\documentclass{./plantilla/umss}

\usepackage[dvipsnames]{xcolor}
\definecolor{bg}{rgb}{0.95,0.95,0.95}

\usepackage{booktabs}
\usepackage{array}
\usepackage{multirow}
\usepackage{enumitem}
\usepackage{minted}
\usepackage{listings}
\usepackage{pdfpages}
\usepackage{graphicx}
\usepackage{url}

\setminted{
  fontsize=\small,
  linenos=true,
  breaklines=true,
  frame=lines,
  framesep=2mm,
  bgcolor=gray!10
}

\setcounter{secnumdepth}{3}

\title{Modelo de Negocio - Sistema Alumni}

\begin{document}

\sujet{Modelo de Negocio - Sistema de Gestión Alumni}
\enseignant{Ing. Rosemary \textsc{Torrico}}
\eleves{\begin{small}
  \begin{itemize}
    \item Azurduy Rivera Pablo Adrian
    \item Bustillos Calderon Jhosua Alejandro
    \item Choquerive Toroya Alfredo Ronald
    \item Lavayen Iguay Adriana
    \item Mercado Apaza Valentina
    \item Prieto Ayala Franco
    \item Rodas Cornejo Nicole Geraldine
    \item Soto Jiménez Angelica Adriana
    \item Terán Mustafá Shamir Leonardo
    \item Vargas Mollo Brenda
  \end{itemize}
\end{small}}

\fairemarges
\fairepagedegarde

\section{Situación Actual}

La relación tradicional entre la universidad pública y sus graduados ha sido un vínculo terminal que cesa tras la titulación. En la Universidad Mayor de San Simón, el contacto se limita a trámites administrativos, careciendo de mecanismos para el seguimiento continuo de la trayectoria laboral. Esto genera un vacío crítico en indicadores de pertinencia curricular e inserción laboral, los cuales son indispensables para la acreditación universitaria.

La infraestructura tecnológica, sostenida por la Unidad de Provisión de Sistemas de Información mediante el sistema webSISS, pierde el rastro del estudiante tras su graduación. El expediente pasa al Departamento de Registros e Inscripciones y al sistema SITRA para trámites notariales, sin integración con plataformas de fidelización. La Dirección Universitaria de Evaluación y Acreditación depende de iniciativas manuales facultativas, lo que deriva en vacíos analíticos.

Esfuerzos aislados en facultades como Ciencias Económicas, Ciencias y Tecnología, y Arquitectura y Ciencias del Hábitat utilizan formularios estáticos o bases de datos desconectadas, impidiendo un seguimiento histórico. Además, se detecta obsolescencia en formatos digitales y ausencia de datos sobre empleabilidad, salarios, motivaciones de ingreso y brecha de género, elementos exigidos en la Dimensión 3 de la acreditación.

\section{Alternativas}

Existen plataformas internacionales que abordan la vinculación laboral con distintos enfoques. Handshake en Estados Unidos conecta estudiantes y empleadores mediante emparejamiento por habilidades. Rikunabi y Mynavi en Japón sostienen el proceso nacional de búsqueda de empleo conocido como shūkatsu mediante fichas de postulación estandarizadas. El sistema SIISEEE en México articula bases de datos dispersas para la toma de decisiones institucionales. JISSEN ALMUNAE en Japón une a estudiantes y egresados desde el primer año con control de acceso segmentado.

\begin{table}[htbp]
  \centering
  \scriptsize
  \renewcommand{\arraystretch}{1.35}
  \resizebox{\textwidth}{!}{%
  \begin{tabular}{>{\raggedright\arraybackslash}p{0.14\textwidth}
                  >{\raggedright\arraybackslash}p{0.21\textwidth}
                  >{\raggedright\arraybackslash}p{0.21\textwidth}
                  >{\raggedright\arraybackslash}p{0.21\textwidth}
                  >{\raggedright\arraybackslash}p{0.21\textwidth}}
    \toprule
    \textbf{Aspecto} & \textbf{Handshake en Estados Unidos} & \textbf{Rikunabi y Mynavi en Japón} & \textbf{SIISEEE en México} & \textbf{JISSEN ALMUNAE en Japón} \\
    \midrule
    Enfoque principal
    & Conectar estudiantes y recién egresados con empleadores para prácticas y primer empleo
    & Sostener el proceso nacional y simultáneo de postulación entre universidades y empresas
    & Integrar información dispersa de estudiantes, egresados y empleadores para decisiones institucionales
    & Vincular a estudiantes en curso con egresados para formación de carrera y apoyo al empleo \\
    \addlinespace
    Mecanismo diferenciador
    & Curaduría y emparejamiento por habilidades y organizaciones estudiantiles
    & Ficha de postulación estandarizada reutilizable y ferias virtuales masivas
    & Articulación de sistemas y bases de datos ya existentes en la universidad
    & Perfiles por tipo de usuario con control de acceso y autenticación por correo o identificación \\
    \addlinespace
    Solución
    & Plataforma dedicada con emparejamiento por habilidades y conexión directa con empleadores verificados
    & Centralización de ofertas y ferias de empleo en un portal único con postulación reutilizable
    & Articulación de instrumentos existentes en un solo sistema para reportes rutinarios y estratégicos
    & Unión de estudiantes y egresados desde el primer año con acceso segmentado por rol \\
    \bottomrule
  \end{tabular}%
  }
  \caption{Comparativa de sistemas de vinculación laboral y trazabilidad de comunidad universitaria.}
  \label{tab:benchmark_internacional}
\end{table}

\section{Perspectiva del Usuario}

Los problemas principales incluyen la pérdida del rastro de los profesionales, la dificultad para levantar métricas de empleabilidad, la falta de intermediación laboral formal, el aislamiento de los estudiantes frente a la industria y la desconexión con la oferta de posgrado. Las causas radican en directorios estáticos, encuestas manuales con baja respuesta y dependencia de redes informales. Las soluciones propuestas abarcan la autenticación y actualización por el propio graduado, reportes en tiempo real, gestión centralizada de postulaciones, emparejamiento con mentores y acceso permanente a recursos universitarios.

\begin{table}[htbp]
  \centering
  \small
  \renewcommand{\arraystretch}{1.35}
  \begin{tabular}{>{\raggedright\arraybackslash}p{0.30\textwidth} >{\raggedright\arraybackslash}p{0.32\textwidth} >{\raggedright\arraybackslash}p{0.32\textwidth}}
    \toprule
    \textbf{Problema} & \textbf{Causa o Impacto Real} & \textbf{Solución} \\
    \midrule
    Pérdida total del rastro de los profesionales tras la titulación. 
    & Imposibilidad de contactar a graduados y bases de datos en hojas de cálculo que caducan en meses. 
    & Autenticación y actualización por el propio graduado y sincronización con redes profesionales externas. \\
    \addlinespace
    Dificultad severa para levantar métricas de empleabilidad e impacto para el CEUB o agencias internacionales como ARCU-SUR. 
    & Encuestas manuales con baja tasa de respuesta y riesgo en auditorías de calidad. 
    & Reportes en tiempo real de absorción laboral, salarios medios, áreas de desempeño y pertinencia curricular. \\
    \addlinespace
    Falta de un canal formal e institucionalizado de intermediación laboral entre la universidad y el sector productivo. 
    & Dependencia de redes informales y ausencia de un canal directo para que las empresas capten talento universitario. 
    & Canales formales de áreas de especialidad y gestión centralizada de postulaciones y pasantías. \\
    \addlinespace
    Aislamiento de los estudiantes de últimos semestres frente a la realidad y exigencias de la industria. 
    & Inseguridad vocacional y demoras prolongadas en la inserción laboral. 
    & Emparejamiento con egresados consolidados, definición de metas y métricas de avance. \\
    \addlinespace
    Desconexión entre la oferta de posgrado o capacitación y los intereses reales de los graduados. 
    & Baja retención de egresados en programas de posgrado propios y pérdida del sentido de pertenencia institucional. 
    & Convocatorias a eventos técnicos, reuniones de cohorte, aranceles preferenciales y acceso permanente a repositorios. \\
    \bottomrule
  \end{tabular}
  \caption{Matriz de alineación problemática y solución del Sistema de Gestión Alumni.}
  \label{tab:problemas_soluciones}
\end{table}

\section{Situación Actual del Seguimiento a la Comunidad Universitaria}

La Dimensión 3 es la única de las cuatro dimensiones de acreditación sin información organizada. No existe seguimiento de egresados, no se conoce la razón de la desvincación de los docentes y no hay contacto regular con las empresas. A continuación se detalla el estado del arte y las propuestas para los actores clave de esta dimensión.

\subsection{Egresados titulados}
En otras universidades se implementan mecanismos robustos de verificación. La Universidad Carlos III de Madrid utiliza dos sistemas separados, uno para mantener el contacto y otro que mide el empleo cruzando la lista de egresados contra el registro estatal de empleo formal, funcionando como una verificación y no como una encuesta. La Universidad de Oxford está obligada por ley a mantener actualizado el contacto de cada egresado para que un organismo independiente lo busque quince meses después de graduarse. La Universidad Autónoma de Madrid utiliza un sistema que completa automáticamente los datos del egresado, aunque la participación se mantiene baja con solo 1200 inscritos de un total de 170000 personas tituladas.

Para la Universidad Mayor de San Simón se propone comenzar cruzando el registro de titulación con el kárdex de notas y la ficha de inscripción, sin esperar a tener acceso a una base de datos externa. El sistema debe completar los datos del egresado de forma automática. Se debe asumir desde el diseño que pocas personas participarán al principio, resolviéndolo con incentivos reales y no solo con buena tecnología.

\subsection{Egresados no titulados}
No existen encuestas específicas para este caso en las universidades consultadas, existiendo únicamente plantillas públicas de encuesta de salida universitaria pensadas para cualquier estudiante que deja la institución. 

Se propone utilizar dicha plantilla agregando una pregunta directa sobre la razón por la cual no se cerró el trámite de titulación, formulada para entender la causa sin presionar a la persona. Además, se propone abrir un expediente por cada caso con etapas simples que van desde materias completas, pasando por el contacto con la persona y el registro del motivo, hasta la toma de una decisión sobre retomar el trámite.

\subsection{Empresas y empleadores}
En Australia existe una encuesta en dos pasos aplicada anualmente desde el año 2016. Primero se encuesta al egresado que ya tiene trabajo solicitando el contacto de su jefe directo, y posteriormente se encuesta a ese jefe preguntando sobre la preparación del egresado en habilidades técnicas y laborales.

Se propone copiar esta idea central, utilizando al egresado como puente para llegar al empleador real. Se deben utilizar también redes sociales y grupos de egresados como canal de contacto, además del correo institucional. Como mejora, se propone volver a encuestar al mismo empleador cada dos años si el egresado sigue trabajando ahí, para observar cambios en su opinión con el tiempo.

\subsection{Docentes}
Un consorcio de universidades diseñó una encuesta única aplicada a dos grupos que luego se comparan, enfocándose en docentes que recibieron una oferta de otra institución y se fueron, y aquellos que recibieron una oferta y decidieron quedarse. La clave es preguntar en el momento en que la persona decide irse, antes de perder el contacto.

Se propone implementar una encuesta corta y confidencial que se envíe automáticamente apenas alguien presenta su renuncia. Como mejora, se debe comparar también con docentes que nunca recibieron ninguna oferta externa, para distinguir entre quienes se fueron por una mejor oferta y quienes lo hicieron por otros motivos como jubilación o insatisfacción. Esta información debe guardarse en el mismo sistema de expedientes por etapas.

\subsection{Estudiantes en curso, postulantes y autoridades}
No se encontraron casos documentados para estos tres actores en las universidades consultadas. 

Para los postulantes se propone abrir un expediente por cada persona con etapas que registren la primera postulación, el resultado y las posteriores postulaciones, permitiendo responder cuántas veces alguien tuvo que intentar entrar y por qué. Para los estudiantes en curso no se necesita un mecanismo nuevo, ya que se les puede encuestar directamente, asegurando que el nuevo sistema consulte los datos que ya existen en el sistema de inscripción de materias. Para las autoridades y consejos, una encuesta directa anual es suficiente.

\subsection{Verificación de datos e incentivos}
En ningún caso real es la propia universidad la que confirma si sus datos de empleo son correctos, siempre intervienen terceros como la Seguridad Social en España, organismos independientes en Inglaterra o centros de investigación contratados por el gobierno en Australia. 

Se propone que el organismo acreditador del Mercosur cumpla parte de ese rol, construyendo el sistema pensando en ese requisito externo. A futuro, se puede plantear que una entidad boliviana externa a la universidad reciba y verifique estos datos de forma independiente.

Para lograr que la encuesta sea una obligación y no una opción, se propone vincularla a un trámite que la persona necesita, como el proceso para recoger el título en físico o pedir certificados. Se debe exigir el dato de contacto actualizado como requisito de egreso y utilizar los beneficios del sistema como incentivo, como condición para seguir usando la bolsa de trabajo o las mentorías.

\section{Situación Legal}

El proceso de acreditación ante el organismo ARCU-SUR, perteneciente al ámbito del Mercosur, establece el marco legal e institucional que exige la demostración de cuatro dimensiones. A diferencia de acreditaciones anteriores realizadas ante el CEUB, este nuevo proceso requiere evidencias documentales y del trabajo real para validar la carrera ante la sociedad. 

La Dimensión 3, correspondiente a la Comunidad Universitaria, es la más débil y la única para la que prácticamente no existe información sistematizada. La normativa exige demostrar que quien estudia informática recibe efectivamente el perfil prometido, que el título sirve para ejercer la profesión, que los docentes son idóneos y que los programas cumplen con los estándares. El incumplimiento de estos parámetros puede resultar en el rechazo de la acreditación, lo que obliga a la universidad a establecer mecanismos legales y formales de seguimiento y registro para validar su oferta académica.

\section{Detalles de los Sistemas Alumni para Solucionar el Problema}

El Programa Alumni de la Universidad Complutense de Madrid sirve como referente. Establece criterios de elegibilidad inclusivos que abarcan desde titulados oficiales hasta estudiantes de movilidad. Define una estructura de membresía jerárquica con categorías base, Plus, Mentor, Honorífico, Ilustre y de Honor. Los derechos incluyen acceso a redes sociales, actividades, servicios universitarios, biblioteca y descuentos, escalando según la categoría alcanzada.

\begin{table}[H]
    \centering
    \begin{tabular}{p{3.5cm}p{9.5cm}}
        \toprule
        \textbf{Categoría} & \textbf{Descripción} \\
        \midrule
        Alumni UCM & Categoría base accesible a todos los perfiles descritos en la sección de elegibilidad. \\
        Alumni UCM Plus+ & Reservada para quienes obtuvieron un título oficial y otorga beneficios adicionales. \\
        Alumni UCM Mentor & Egresados con experiencia profesional relevante que asesoran a otros de forma voluntaria y son designados por las facultades. \\
        Alumni UCM Honorífico & Personas físicas o jurídicas destacadas por su apoyo a la universidad, designadas por el Comité de Dirección. \\
        Alumni UCM Ilustre & Egresados de prestigio reconocido en su ámbito profesional o de conocimiento. \\
        Alumni UCM de Honor & Personalidades de máxima relevancia social formadas en la Complutense. \\
        \bottomrule
    \end{tabular}
    \caption{Categorías de membresía del programa Alumni UCM.}
    \label{tab:categorias-alumni-ucm}
\end{table}

Además del modelo de membresía, se proponen funciones útiles basadas en sistemas de referencia como el de la universidad chilena Alumni UC. Se debe priorizar la bolsa de trabajo y las mentorías, ya que cada una sirve para conectar a dos actores de la Dimensión 3 con una sola función. Se implementará un catálogo simple de beneficios propios, evitando descuentos comerciales poco realistas. Se utilizarán los certificados y trámites en línea como la razón práctica para que el egresado vuelva al sistema, en lugar de depender de una encuesta anual aislada.

\subsection{Sistematización mediante CiviCRM}
Varios de los casos citados en España utilizan el mismo programa de código abierto llamado CiviCRM, diseñado para asociaciones y organizaciones sin fines de lucro. No es una sola herramienta, sino un registro central de contactos al que se le pueden activar o desactivar módulos según lo que la organización necesite, como membresías con cuota, eventos con cobro, correo masivo, expedientes por etapas y reportes. Esto explica por qué algunas universidades cobran una cuota únicamente a quien decide afiliarse formalmente a la asociación.

\section{Actores}

El sistema contempla actores que interactúan con la plataforma de acuerdo con sus necesidades y funciones dentro de la comunidad de egresados, estableciendo la vinculación entre estos, la carrera y el mercado laboral.

\subsection{Egresado}
El egresado es la persona que ha concluido su formación académica y utiliza la plataforma para mantener y fortalecer su vínculo con la universidad y el entorno profesional. Este actor gestiona su perfil profesional, registra información relacionada con su experiencia, consulta oportunidades laborales, participa en procesos de mentoría y responde encuestas sobre su situación profesional.

\subsection{Empresa o Reclutador}
La empresa o reclutador es el actor externo que utiliza la plataforma para establecer contacto con profesionales egresados y acceder a oportunidades de vinculación laboral. Este actor publica ofertas de empleo y busca perfiles de egresados de acuerdo con las habilidades y áreas profesionales requeridas.

\subsection{Administrador de la Carrera}
El administrador de la carrera es el actor encargado de gestionar y supervisar la información institucional relacionada con los egresados. Accede a información consolidada sobre la comunidad, gestiona encuestas y consulta indicadores relacionados con la empleabilidad y trayectoria profesional.

\subsection{Mentor}
El mentor es un egresado con experiencia profesional que participa voluntariamente en actividades de orientación y acompañamiento dirigidas a otros egresados. Este actor ofrece espacios de mentoría y gestiona la disponibilidad de horarios para las sesiones de orientación profesional.

\section{Referencias}

\begin{itemize}
    \item \textbf{Universidad Carlos III de Madrid:} Ficha de carrera en Gradopedia con cita a registros de afiliación a la Seguridad Social (\url{https://gradopedia.es/grados/ciencia-e-ingenieria-de-datos/uc3m-madrid}), página oficial de empleabilidad (\url{https://www.uc3m.es/estudios/empleo-practicas}) e informe del Ministerio de Educación de España sobre inserción laboral.
    \item \textbf{Universidad de Oxford y Agencia de Estadísticas de Educación Superior:} Servicio de Carreras de Oxford (\url{https://www.careers.ox.ac.uk/the-graduate-outcomes-survey}) y Office for Students, organismo regulador inglés (\url{https://www.officeforstudents.org.uk/for-providers/quality-and-standards/graduate-outcomes-survey/}).
    \item \textbf{Universidad Autónoma de Madrid:} Artículo sobre el desarrollo de su plataforma AlumniUAM (\url{https://www.ixiam.com/es/blog/la-uam-confia-en-ixiam-el-desarrollo-de-su-plataforma-alumniuam/}).
    \item \textbf{Plantilla de encuesta de salida universitaria:} QuestionPro (\url{https://www.questionpro.com/survey-templates/university-exit/}).
    \item \textbf{Australia, encuesta de satisfacción del empleador:} Descripción oficial del mecanismo (\url{https://srcentre.com.au/project/employer-satisfaction-survey/}) e informe nacional 2021 (\url{https://www.qilt.edu.au/docs/default-source/default-document-library/2021-ess-national-report.pdf}).
    \item \textbf{Consorcio universitario, encuesta de retención y salida de docentes:} Página oficial de la Universidad de Harvard (\url{https://coache.gse.harvard.edu/faculty-retention-and-exit-survey}).
    \item \textbf{Sistema Alumni de referencia en Chile:} Beneficios y credencial virtual (\url{https://alumni.uc.cl/actualidad/vuelve-a-vivir-la-universidad-en-mi-portal-uc/}), bolsa de trabajo (\url{https://alumni.uc.cl/beneficio/mercado-laboral-uc/}), mentorías (\url{https://mentoriaslaborales.uc.cl/programa-de-mentorias}) y página principal (\url{https://alumni.uc.cl/}).
    \item \textbf{CiviCRM:} Wikipedia, lista de módulos (\url{https://es.wikipedia.org/wiki/CiviCRM}).
\end{itemize}

\end{document}